%%%%%%%%%%%%%%%%%%%%%%%%%%%%%%%%%%%%%%%%%%%%%%%%%%%%%
%			GENERÁTOR v. 18.09 				%
%%%%%%%%%%%%%%%%%%%%%%%%%%%%%%%%%%%%%%%%%%%%%%%%%%%%%
% Tento soubor slouží k jednotlivé zkusné kompilaci 
% písniček či k kompilaci samostatných pdf souborů 
% pomocí bashového skriptu kompilace.sh ve stejné 
% složce. Z důvodu hromadné kompilace musí poslední 
% dva řádky tohoto souboru zůstat nezměněny.
%%%%%%%%%%%%%%%%%%%%%%%%%%%%%%%%%%%%%%%%%%%%%%%%%%%%%
%			Jak kompilovat jednotlivé písně?        %
%%%%%%%%%%%%%%%%%%%%%%%%%%%%%%%%%%%%%%%%%%%%%%%%%%%%%
%	1. Do míst PÍSNIČKA upravte následnující řádek, 
%	   aby vypadal takto: 
%	   \input{../songy/[JMÉNO SOUBORU KOMPILOVNÉ PÍSNIČKY].tex}
%	2. Soubor písničky musí být přeložitelný a musí 
%	   se nacházet ve složce ../songy.
%%%%%%%%%%%%%%%%%%%%%%%%%%%%%%%%%%%%%%%%%%%%%%%%%%%%%
%			Jak psát soubory songů?                 %
%%%%%%%%%%%%%%%%%%%%%%%%%%%%%%%%%%%%%%%%%%%%%%%%%%%%%
%	1. Více návodu je k tomuto napsáno v souboru 
%      ../songy/00Songtemplate. 
%%%%%%%%%%%%%%%%%%%%%%%%%%%%%%%%%%%%%%%%%%%%%%%%%%%%%
%			Jak kompilovat celý zpěvník?			%
%%%%%%%%%%%%%%%%%%%%%%%%%%%%%%%%%%%%%%%%%%%%%%%%%%%%%
%	1. Více návodu je k tomuto napsáno v souboru
%	   ../Cely_zpevnik/zpevnik.tex.
%%%%%%%%%%%%%%%%%%%%%%%%%%%%%%%%%%%%%%%%%%%%%%%%%%%%%

% Hlavička START
\documentclass[openany,12pt]{memoir}
\usepackage[utf8]{inputenc} 
\usepackage[czech]{babel}
\usepackage[T1]{fontenc}
\usepackage[top=1.5cm, bottom=2cm, left=2cm, right=2cm]{geometry}  % --> NASTAVENÍ OKRAJŮ
\usepackage{fancyhdr}
\usepackage{graphicx}
\usepackage{xwatermark}
\usepackage{xcolor}
\usepackage{changepage}
\usepackage{pdfpages}
\usepackage{lettrine}
\usepackage{indentfirst}  %Důležité pro formátování

%%%%%%%%%%%%%%%%%%%%%%%%%%%%%%%%%%%%%%
%  FONT                              %
%%%%%%%%%%%%%%%%%%%%%%%%%%%%%%%%%%%%%%
\usepackage{amssymb}
\usepackage{tgschola}



%%%%%% Package na zpěvník
\usepackage[full]{leadsheets}%http://mirrors.nic.cz/tex-archive/macros/latex/contrib/leadsheets/leadsheets_en.pdf   --> dokumentace	
\definesongtitletemplate{empty}{} 
\setchords{
format = \bfseries \sffamily,   %tučné akordy
minor = {mi},% 
input-notation = {german},%
output-notation = {german}%
}
\definesongtitletemplate{empty}{} 


\newlength{\drop}
% VODOZNAK
\newwatermark[pages=1-,color=red!50,angle=0,scale=2, xpos=0,ypos=0]{\includegraphics[width=5cm]{obr/pozadi2.jpg}} %--> dvojka na pozadí


%%%%%%%%%%%%%%%%%%%%%%%%%%%%%%%%%%%%%%%%%%%%%%%%%%
%		 Vlastní příkazy
\newcounter{Slokočet}   %Automatické číslování slok
\newcommand{\mezera}{
\phantom{.}

}   %Horizontální odsazení slok (poněkud blbě zadefinovaný, ale jinak se formát rozbije jako wtf prostě)
\newcommand{\stred}{5.2cm}   %%% Na zarovnání slok doprostřed, pozn. automatičtější zarovnávání na střed nejde
\newcommand{\carka}{,\:}
\newcommand{\m}[1]{\color{white}{#1}}  %Pro akordy
\newcommand{\ap}{'}	%Pro apostrof
\newcommand{\elipsa}{\kern\fontdimen3\font} %Příkaz pro lepší zacházení s výpustkami (=...); je to vpodstatě jen mezera mezi tečkama výpustky
\newcommand{\pindent}{17.62482 pt} %Správná velikost \parindentu u layoutu se dvěma minipageama
\newcommand{\predtitle}{\huge}
\newcommand{\mezisloupci}{\phantom{TT}} %Místo mezi dvěma sloupci na jedné stránce
\newcommand{\z}{\hspace*{\fill}\null}

%%% Možné velikosti písem 
\newcommand{\normalni}{\normalsize}
\newcommand{\velky}{\fontsize{14.4}{15}\selectfont}
\newcommand{\vetsi}{\fontsize{15}{16}\selectfont}
\newcommand{\nejvetsi}{\fontsize{16}{17}\selectfont}
\newcommand{\nejnejvetsi}{\fontsize{17}{19}\selectfont}

%%% Stará definice sloky spoléhající na indenty
%\newlength{\pismeno}
%\settowidth{\pismeno}{x} %Tohle není moc ideální velikost, ale funguje
%\newif\ifslokavelka
%\slokavelkafalse
%\newcommand{\sloka}{
%\ifnum \value{Slokočet}>8  %Pokud je sloka dvouciferná
%\mezera \noindent \addtocounter{Slokočet}{1} \hspace*{-\pismeno}\arabic{Slokočet}.
%\else %Pokud jen jednociferná
%\mezera \noindent \addtocounter{Slokočet}{1} \arabic{Slokočet}. 
%\fi
%} 	%sloka, která se automaticky čísluje

\newcommand{\distanc}{\:}  %Vzdálenost čísla sloky před slokou
\newlength{\delkaargumentu}
%%% Sloka s automatickým číslováním
\newcommand{\sloka}{%
\addtocounter{Slokočet}{1}% Zvýší se o 1 počet slok
\mezera%  Sloka se odsadí vertikálně
\settowidth{\delkaargumentu}{\arabic{Slokočet}.\distanc}% Zde se určí délka odsazení 
\hspace*{-\delkaargumentu}%
\arabic{Slokočet}.\distanc%
\ignorespaces% Aby nevznikaly zbytečné mezery
}


%%% Sloka s vlastním argumentem
\newcommand{\ssloka}[1]{%     
\settowidth{\delkaargumentu}{#1\distanc}
\mezera%
\hspace*{-\delkaargumentu}%
#1\distanc%
\ignorespaces%
}  

%%% Refrén
\newcommand{\refren}[1][0]{%  Nepovinný argument sděluje, kolikátý refrén toto je, bez argumentu se vytiskne pouze refrén
\ifnum #1>0 %Pokud nepovinný argument existuje
\mezera%
\settowidth{\delkaargumentu}{\textbf{R$_{\text{#1}}$:}\distanc}%
\hspace*{-\delkaargumentu}%
\textbf{R$_{\text{#1}}$:}\distanc%
\ignorespaces%
\else %Pokud nepovinný argument neexistuje
\mezera%
\settowidth{\delkaargumentu}{\textbf{R:}\distanc}%
\hspace*{-\delkaargumentu}%
\textbf{R:}\distanc%
\ignorespaces%
\fi
}

\newcommand{\predehra}{\ssloka{\textbf{Předehra:}}}

%%% Capo
\newcommand{\kapodastr}[1]{%
\textit{Capo {#1}}
}


\addto\captionsczech{\renewcommand{\contentsname}{Seznam písní}}

%%%%%%%%%%%%%%%%%%%%%%%%%%%%%%%%%%%%
%    FORMÁTOVÁNÍ                   %
%%%%%%%%%%%%%%%%%%%%%%%%%%%%%%%%%%%%

%%% Vlevo zarovnaný text s blokem zarovnaným na střed
\usepackage{varwidth}% http://ctan.org/pkg/varwidth
\newenvironment{centerjustified}{%
  \begin{center} % so the minipage is centered
  \begin{varwidth}[t]{\textwidth}	
  \raggedright % so the minipage's text is left justified
  \setlength{\parindent}{\pindent}
}{%
  \end{varwidth}
  \end{center}
}


%Hlavnička END
\begin{document}
\sffamily %Sans-serif pro lepší čitelnost

\pagestyle{empty}
\newgeometry{top=1.5cm, bottom = 0cm, left = 2cm, right = 2cm}
	
\newpage

%Pro kompilaci je důležitý, aby ten \input a \end{document} zůstali na posledních
%dvou řádcích!!!!

\velky
\newpage
% PÍSNIČKA
%\documentclass[../main.tex]{subfiles}

\begin{song}{title=\centering What Makes You Beautiful \\\normalsize One Direction  \vspace*{-0.3cm}}  %% sem se napíše jméno songu a autor
\moveright 3cm \vbox{      %Varianta č. 1  ---> Jeden sloupec zarovnaný na střed	


\sloka
You're ^{\,\,D}insecure ^{G} 

Don't know what ^{A}for 

You're turning ^{D{\color{white}\_}}heads when you ^{G{\color{white}\_}}walk through the ^{A{\color{white}\_}}door 

Don't need make ^{D}up 

^{G}To cover ^{A}up 

Being the ^{D}way that you ^{G}are is ^{A}enough 

%\moveright -0.2cm \vbox{\ssloka \textbf{R\textsubscript{1}:}
^{D{\color{white}\_\_\_}}Everyone ^{G\,\,}else in the ^{A{\color{white}\_}}room can see it 

^{D{\color{white}\_\_\_}}Everyone ^{G}else but ^{A}you
 
%\moveright -0.2cm \vbox{\ssloka \textbf{R\textsubscript{2}:}
Baby you ^{D{\color{white}\_}}light up my ^{G{\color{white}\_}}world like ^{A}nobody else 

The way that ^{D}you flip your ^{G}hair gets me ^{A{\color{white}\_\_\_\_}}overwhelmed 

But you when ^{D{\color{white}\_}}smile at the ^{G{\color{white}\_\_}}ground it ain't ^{A}hard to tell 

You don't ^{\,\,\,D (Hmi)}know oh ^{G}oh

^{A}You don't know you're ^{{\color{white}\_\_}D}beautiful 

%\moveright -0.2cm \vbox{\ssloka \textbf{R\textsubscript{3}:}
If only ^{A}you saw what ^{H}I can see 

You'll ^{D{\color{white}\_\_\_\_}}understand why I ^{G}want you so ^{A{\color{white}\_\_\_\_}}desperately 

Right now I'm ^{D{\color{white}\_\_\_}}looking at ^{G}you and I ^{A}can't believe 

You don't ^{\,\,\,D (Hmi)}know oh ^{G}oh

^{A}You don't know you're ^{{\color{white}\_\_}D}beautiful oh - ^{G}oh  

But ^{A}that's what makes you ^{{\color{white}\_\_}D}beautiful 

\sloka
So c-come on 

You got it wrong

To prove I'm right I put it in a song  

I don't why

You're being shy

And turn away when I look into your eyes 


}
\newpage
\moveright 3cm \vbox{      %Varianta č. 1  ---> Jeden sloupec zarovnaný na střed	

%\moveright -0.2cm \vbox{\ssloka \textbf{R\textsubscript{1}:}
^{D{\color{white}\_\_\_}}Everyone ^{G\,\,}else in the ^{A{\color{white}\_}}room can see it 

^{D{\color{white}\_\_\_}}Everyone ^{G}else but ^{A}you
 
%\moveright -0.2cm \vbox{\ssloka \textbf{R\textsubscript{2}:}
Baby you ^{D{\color{white}\_}}light up my ^{G{\color{white}\_}}world like ^{A}nobody else 

The way that ^{D}you flip your ^{G}hair gets me ^{A{\color{white}\_\_\_\_}}overwhelmed 

But you when ^{D{\color{white}\_}}smile at the ^{G{\color{white}\_\_}}ground it ain't ^{A}hard to tell 

You don't ^{\,\,\,D (Hmi)}know oh ^{G}oh

^{A}You don't know you're ^{{\color{white}\_\_}D}beautiful 

%\moveright -0.2cm \vbox{\ssloka \textbf{R\textsubscript{3}:}
If only ^{A}you saw what ^{H}I can see 

You'll ^{D{\color{white}\_\_\_\_}}understand why I ^{G}want you so ^{A{\color{white}\_\_\_\_}}desperately 

Right now I'm ^{D{\color{white}\_\_\_}}looking at ^{G}you and I ^{A}can't believe 

You don't ^{\,\,\,D (Hmi)}know oh ^{G}oh

^{A}You don't know you're ^{{\color{white}\_\_}D}beautiful oh - ^{G}oh  

But ^{A}that's what makes you ^{{\color{white}\_\_}D}beautiful 

\sloka
^{D}Nana Nana ^{G}Nana Nana  ^{A}Nana Nana na

^{D}Nana Nana ^{G}Nana Nana  ^{A}Nana Nana na 

^{G}Nana Nana^{A} Nana Nana   

%\moveright -0.2cm \vbox{\ssloka \textbf{R\textsubscript{2}:}

%\moveright -0.2cm \vbox{\ssloka \textbf{R\textsubscript{2}:}

%\moveright -0.2cm \vbox{\ssloka \textbf{R\textsubscript{3}:}
 

}
\setcounter{Slokočet}{0}
\end{song}

\end{document}
